\documentclass[11pt]{article}

\usepackage[margin=2.2cm]{geometry}
\usepackage[utf8]{inputenc}
\usepackage[english]{babel}
\usepackage{enumitem}
\usepackage{hyperref}

\begin{document}

\begin{center}
    {\LARGE \textbf{Dr. Eng. Ramiro M. Irastorza}}\\[0.2cm]
    Independent Researcher -- CONICET\\
    Institute of Physics of Liquids and Biological Systems (IFLySiB-CONICET-UNLP)\\
    La Plata, Argentina\\

    \vspace{0.3cm}

    \href{mailto:rirastorza@iflysib.unlp.edu.ar}{rirastorza@iflysib.unlp.edu.ar} \\

    \href{https://mel.iflysib.unlp.edu.ar/}{Grupo MELyMCB-IFLySiB}

    Google Scholar:
    \href{https://scholar.google.com/citations?hl=en&user=0LriuSsAAAAJ&view_op=list_works&sortby=pubdate}{Link}

    ORCID:
    \href{https://orcid.org/0000-0002-6455-3574}{0000-0002-6455-3574}

    GitHub:
    \href{https://github.com/rirastorza}{github.com/rirastorza}
\end{center}

% \vspace{0.5cm}

\section*{Education}

\begin{itemize}[leftmargin=*]
    \item \textbf{Ph.D. in Electronic Engineering}, National University of La Plata (UNLP), Argentina, 2010.

    Dissertation: \textit{System Identification: Application to the In Vitro Evaluation of Bone Quality in Human Trabecular Tissue}.

    \item \textbf{Electronic Engineer}, National University of La Plata (UNLP), Argentina, 2005.
\end{itemize}

\section*{Current Positions}

\begin{itemize}[leftmargin=*]
    \item \textbf{Independent Researcher}, CONICET (2026--Present).
    \item \textbf{Full Professor}, Automation Laboratory, Mechanical Engineering Department, UTN-FRLP (2025--Present).
    \item Vice Director, Center for Experimental and Computational Mechanics (CIMEC), UTN-FRLP.
\end{itemize}

\section*{Research Interests}

\begin{itemize}[leftmargin=*]
    \item Computational bioengineering
    \item Finite element modeling
    \item Microwave tomography and inverse problems
    \item Radiofrequency and microwave medical therapies
    \item Biomedical image reconstruction
    \item Thermal and dielectric properties of biological tissues
    \item Scientific computing and machine learning
\end{itemize}

\section*{Research Funding and Leadership}

\begin{itemize}[leftmargin=*]
    \item Principal Investigator of the UTN project:
    \textit{Biomedical and Biomimetic Problems Solved Using Finite Elements} (2023--2026).

    \item Principal Investigator of the national research project:
    \textit{Microwave Tomography: Reconstruction Algorithms, Experimental Validation and Applications} funded by the National Agency for Scientific and Technological Promotion (2022--2024).

    \item Co-Principal Investigator of the CONICET project:
    \textit{Inverse Problems in Biomedical Engineering Using Sparse Representations and Artificial Intelligence} (2023--2024).
\end{itemize}

\section*{Selected International Research Project Participation}

\begin{itemize}[leftmargin=*]
 \item \textbf{Improvement of Advanced Ablative Therapies through the Control of Tissue and Cellular Behavior Using Electromagnetic Fields}.
Participants: Polytechnic University of Valencia, Neuroscience and Complex Systems Studies Center (CONICET), Institute of Electronics, Signal Processing and Control Research (LEICI-CONICET), and IFLySiB-CONICET.
Funded by the Ministry of Economy and Competitiveness, Spain (2023--2025).

\item \textbf{Radiofrequency- and Microwave-Based Technologies for Minimally Invasive Surgery}.
Participants: Polytechnic University of Valencia and IFLySiB-CONICET.
Funded by the Ministry of Economy and Competitiveness, Spain (2015--2017).

\item \textbf{Mise en place d’un laboratoire international pour la recherche sur les matériaux structures recouvrements pour applications médicales}.
Participants: Université Laval, INTI, and the Biomechanics Laboratory, School of Engineering, University of Buenos Aires.
Funded by the Agence Universitaire de la Francophonie (AUF) (2010--2012).

\end{itemize}


\section*{Graduate Teaching}

\begin{itemize}[leftmargin=*]
    \item \textbf{Finite Element Method with Open-Source Software} (UTN-FRLP)
    \item \textbf{Computational Tools for Scientists} (UNLP / UTN)
    \item \textbf{Statistical Methods} (UNAJ)
    \item \textbf{Applied Mathematics Complements} (UNAJ)
    \item \textbf{Radiofrequency and Biological Tissues} (Polytechnic University of Valencia, Spain)
\end{itemize}

\section*{Human Resources Training}

\begin{itemize}[leftmargin=*]
    \item Supervision and co-supervision of Ph.D. students in biomedical engineering, computational mechanics, and microwave imaging.
    \item Advisor of undergraduate and graduate theses at UTN, UNAJ, and UNLP.
    \item Supervisor of CONICET and national research fellowship recipients.
\end{itemize}

\section*{Selected Publications}

\begin{enumerate}[leftmargin=*]

\item Berjano E., \textbf{Irastorza R.M.},
\textit{Positioning of the dispersive electrode and its effect on the safety and efficacy of radiofrequency ablation},
EP Europace, 2025.

\item Collavini S., Pérez J.J., Berjano E., Fernández-Corazza M., Oddo S., {\bf Irastorza R.M.},
\textit{Impact of surrounding tissue-type and peri-electrode gap in stereoelectroencephalography guided (SEEG) radiofrequency thermocoagulation (RF-TC): a computational study},
International Journal of Hyperthermia, 2024.

\item Cervantes M.J., Basiuk L.O., González-Suárez A., Carlevaro C.M., \textbf{Irastorza R.M.},
\textit{Low-Frequency Electrical Conductivity of Trabecular Bone: Insights from In Silico Modeling},
Mathematics, 2023.

\item \textbf{Irastorza R.M.}, Maher T., Barkagan M., Liubasuskas R., Berjano E., d’Avila A.,
\textit{Anterior vs. posterior position of dispersive patch during radiofrequency catheter ablation: insights from in silico modelling},
EP Europace, 2023.

\item González-Suárez A., Pérez J.J., \textbf{Irastorza R.M.}, D'Avila A., Berjano E.,
\textit{Computer modeling of radiofrequency cardiac ablation: 30 years of bioengineering research},
Computer Methods and Programs in Biomedicine, 2022.

\item \textbf{Irastorza R.M.}, Bovaira M., García-Vitoria C., Muñoz V., Berjano E.,
\textit{Effect of the relative position of electrode and stellate ganglion during thermal radiofrequency ablation},
International Journal of Hyperthermia, 2021.

\item Fajardo J.E., Lotto F.P., Vericat F., Carlevaro C.M., \textbf{Irastorza R.M.},
\textit{Microwave Tomography with Phaseless Data on the Calcaneus by Means of Artificial Neural Networks},
Medical \& Biological Engineering \& Computing, 2019.

\item \textbf{Irastorza R.M.}, d'Avila A., Berjano E.,
\textit{Thermal latency adds to lesion depth after application of high-power short-duration radiofrequency energy},
Journal of Cardiovascular Electrophysiology, 2018.

\item \textbf{Irastorza R.M.}, Trujillo M., Berjano E.,
\textit{How coagulation zone size is underestimated in computer modeling of RF ablation by ignoring the cooling phase},
International Journal for Numerical Methods in Biomedical Engineering, 2017.

\item \textbf{Irastorza R.M.}, Blangino E., Carlevaro C.M., Vericat F.,
\textit{Modeling of the dielectric properties of trabecular bone samples at microwave frequency},
Medical \& Biological Engineering \& Computing, 2014.

\end{enumerate}

\section*{Professional Activities}

\begin{itemize}[leftmargin=*]
    \item Reviewer for international journals including:
    \begin{itemize}
        \item International Journal of Hyperthermia
        \item Medical Engineering \& Physics
        \item Progress In Electromagnetics Research
        \item International Journal for Numerical Methods in Biomedical Engineering
    \end{itemize}

    \item Guest Editor:
    \textit{Biophysical- and Engineering-Oriented Approaches to Energy-Based Cardiac Ablation Techniques},
    Journal of Cardiovascular Development and Disease.

    \item Co-Coordinator of the Information Technology, Communications and Electronics Area, National Agency for Scientific and Technological Promotion FONCyT (PICT calls 2022--2025).

    \item Organizer of PyDay La Plata (2022--2025).
\end{itemize}

\end{document}
